\section{Conclusion}
This study revises a two-stage stock-index volatility-proxy forecasting exercise by separating numerical reproduction, data correction, and predictive evaluation. A sampling-compatible XGBoost build reproduces all 6888 archived first-stage predictions. The original DEC environment cannot be uniquely reconstructed, so the current method uses an explicitly defined profiled objective, analytic gradients, deterministic bounded multistart search, and stationarity checks. The confirmed SSE50 holiday record is removed before feature and target construction, and XGB-AR supplies a residual-only ablation.

The substantive conclusion is conditional on the loss function. Relative to XGB-Market and XGB-AR, revised DEC reduces QLIKE but increases log-MSE in all three markets. MAE improves in the two Chinese markets but worsens in Nasdaq. GARCH remains the best QLIKE benchmark in these comparisons. Improved numerical optimization is valuable for method correctness even when a test metric deteriorates: Nasdaq's validation-selected switch from Beta to exponential Almon improves MSE relative to archived DEC but worsens QLIKE.

Interval evaluation does not support a parallel advantage. On the 1859 common CSI300 dates, the main DEC interval score is 10.1995, worse than HAR's 9.6973, XGB-Market's 9.8534, and XGB-AR's 9.8793. Non-rejection of coverage diagnostics in the main setting does not establish validity, and the 95\%/120-window setting still shows rejection of violation independence and conditional coverage. These are intervals for an observed daily proxy under a specific residual-calibration rule, not confidence intervals for latent variance.

The methodological value is the explicit separation of archived evidence from revised findings, the loss-dependent residual ablation, and the numerical checks. The study does not establish a new general forecasting theorem, an exhaustive advantage over modern alternatives, a global optimizer, or a causal policy mechanism. Weight parameters at their bounds should not be treated as robust economic transmission estimates.

Important limitations remain. The daily proxy is noisy, the log transformation changes the prediction functional, real-time EPU release and revision histories are missing, and the full upstream calendar and data lineage are not independently audited. Historical results were already observed, and the comparison markets are not untouched holdouts. HAC and Holm calculations are therefore exploratory rather than confirmatory. Further work should prioritize a locked design, genuinely available-at-origin inputs, an independent future holdout, and stronger directly implemented benchmarks before making operational forecasting claims.
